% Clase 2 — Búsqueda Desinformada (parte 1)
% Lectura: AIMA Secciones 3.1–3.3
\section{Búsqueda Desinformada I: Formulación del Problema}

\subsection{Formulación de un Problema de Búsqueda}

\begin{definition}[Problema de búsqueda]
Un problema de búsqueda se define con cinco componentes:
\begin{enumerate}
    \item \textbf{Estado inicial}: $s_0$
    \item \textbf{Modelo de transición}: $\text{RESULT}(s, a) \rightarrow s'$
    \item \textbf{Función de costo de acción}: $c(s, a, s')$
    \item \textbf{Test de meta}: $\text{IS-GOAL}(s) \rightarrow \{T, F\}$
    \item \textbf{Espacio de estados}: grafo de todos los estados alcanzables
\end{enumerate}
\end{definition}

La \textbf{solución} es una secuencia de acciones que lleva del estado inicial a un estado meta. Una \textbf{solución óptima} minimiza el costo total del camino.

\subsection{Nodos vs. Estados}

Un \textbf{nodo} en el árbol de búsqueda contiene:
\begin{itemize}
    \item El estado que representa
    \item El nodo padre
    \item La acción que generó el nodo
    \item El costo del camino $g(n)$
    \item La profundidad
\end{itemize}

\subsection{Medidas de Desempeño de Algoritmos}

\begin{tabular}{ll}
\toprule
\textbf{Criterio} & \textbf{Descripción} \\
\midrule
Completitud & ¿Encuentra solución si existe? \\
Optimalidad & ¿Encuentra la solución de menor costo? \\
Complejidad en tiempo & Nodos generados/expandidos \\
Complejidad en espacio & Nodos en memoria simultáneamente \\
\bottomrule
\end{tabular}

Los parámetros usuales son $b$ (factor de ramificación), $d$ (profundidad de la solución más superficial) y $m$ (profundidad máxima del árbol).
